\section*{Chapitre \ref{ch:b1:organic-redox} --- Oxydation et réduction en chimie organique}

\begin{solution}{exo:b1:organic-redox:1}
Éthanol~: \ce{CH3} $-3$, \ce{CH2OH} $-1$. Éthanal~: \ce{CH3} $-3$, CHO $+1$.
Acide éthanoïque~: \ce{CH3} $-3$, COOH $+3$. Propanone~: \ce{CH3} $-3$
(deux fois),
C=O $+2$. Acide méthanoïque~: $+2$.
\end{solution}

\begin{solution}{exo:b1:organic-redox:2}
Propan-1-ol~: propanal, puis acide propanoïque. Propan-2-ol~: propanone.
2-Méthylpropan-2-ol~: pas de réaction (la couleur orange persiste).
\end{solution}

\begin{solution}{exo:b1:organic-redox:3}
\ce{CH3COCH3 + 2H+ + 2e- -> CH3CH(OH)CH3}~;
\ce{CH3CHO + 2H+ + 2e- -> CH3CH2OH}.
\end{solution}

\begin{solution}{exo:b1:organic-redox:4}
\ce{NaBH4}~: butan-1-ol, butan-2-ol, pas de réaction, pas de réaction
(l'acide est seulement déprotoné). \ce{LiAlH4}~: butan-1-ol, butan-2-ol,
butan-1-ol et éthanol, butan-1-ol.
\end{solution}

\begin{solution}{exo:b1:organic-redox:5}
\ce{5CH3CH(OH)CH3 + 2MnO4- + 6H+ -> 5CH3COCH3 + 2Mn^2+ + 8H2O}.
$n = 1.20/60.0 = \qty{0.0200}{mol}$~; permanganate $\frac25 \times 0.0200 =
\qty{8.0e-3}{mol}$, soit \qty{400}{mL} de la solution à \qty{0.020}{mol/L}.
\end{solution}

\begin{solution}{exo:b1:organic-redox:6}
Chauffer l'alcool avec l'\omterm{def:b1:nernst:couple}{oxydant} dans un ballon monté pour une
distillation, à une température comprise entre les deux températures
d'ébullition~: le propanal (\qty{48}{\celsius}) distille au fur et à
mesure de sa formation et échappe à une oxydation ultérieure~; le
propan-1-ol (\qty{97}{\celsius}) reste dans le ballon.
\end{solution}

\begin{solution}{exo:b1:organic-redox:7}
Un doublet B--H de \ce{BH4-} attaque le carbone du carbonyle, le
doublet $\pi$ de C=O passant sur l'oxygène~; l'\omterm{def:b1:alcohol-activation:alkoxide}{alcoolate} capte un proton
de l'éthanol. L'hydrogène porté par le carbone du produit
(\ce{CH3CH2OH}, l'un des deux de C1) provient du réactif~; celui de
l'oxygène provient du solvant.
\end{solution}

\begin{solution}{exo:b1:organic-redox:8}
Le butan-2-ol est \omterm{def:b1:stereochemistry-in-depth:stereocentre}{chiral}, mais la cétone plane est attaquée également
sur ses deux faces~: on obtient un \omterm{def:b1:stereochemistry-in-depth:enantiomers}{mélange racémique}, optiquement
inactif.
\end{solution}

\begin{solution}{exo:b1:organic-redox:9}
\ce{3CH3CH2OH + 2Cr2O7^2- + 16H+ -> 3CH3COOH + 4Cr^3+ + 11H2O}~: le
dichromate orange (chrome $+$VI) devient du chrome(III) vert.
\end{solution}

\begin{solution}{exo:b1:organic-redox:10}
2-Méthylbutan-2-ol~: l'éthanal + \ce{CH3CH2MgBr} donne le butan-2-ol
(une addition, sans oxydoréduction)~; oxydation en butanone~;
\ce{CH3MgBr} sur la butanone, puis hydrolyse. Une oxydation.
Pentan-3-ol~: propanal + \ce{CH3CH2MgBr} en une étape, sans
oxydoréduction.
\end{solution}

\begin{solution}{exo:b1:organic-redox:11}
1. Protéger la cétone sous forme d'\omterm{def:b1:acetals-protection:hemiacetal}{acétal} cyclique (éthane-1,2-diol,
acide, Dean--Stark). 2. \ce{LiAlH4} dans l'éther anhydre réduit l'ester
en alcool primaire (et méthanol). 3. Une solution aqueuse acide retire
l'\omterm{def:b1:acetals-protection:hemiacetal}{acétal}~: 5-hydroxypentan-2-one, \ce{CH3CO(CH2)3OH}.
\end{solution}

\begin{solution}{exo:b1:organic-redox:12}
\ce{RCH2OH}~: carbone fonctionnel à $-1$~; \ce{RCOOH}~: $+3$~; quatre
électrons, comme dans
\ce{RCH2OH + H2O -> RCOOH + 4H+ + 4e-}. Du méthane ($-4$) au dioxyde de
carbone ($+4$)~: huit électrons, \ce{CH4 + 2H2O -> CO2 + 8H+ + 8e-}.
\end{solution}

\begin{solution}{pb:b1:organic-redox:1}
\textbf{1.} C1 (qui porte OH)~: 0~; les cinq \ce{CH2}~: $-2$.
\textbf{2.} C1 (C=O)~: $+2$~; les cinq \ce{CH2}~: $-2$.
\textbf{3.} Les deux carbones de COOH~: $+3$~; les quatre \ce{CH2}~: $-2$.
\textbf{4.} C1 ($0 \to +3$) et son voisin C6 ($-2 \to +3$), qui devient
le second carbone carboxylique~; la liaison C1--C6 est rompue.
\textbf{5.} $3 + 5 = 8$ électrons.
\textbf{6.} L'acide nitrique est assez fort pour rompre la liaison C--C
voisine du carbonyle, ce que ne font pas des \omterm{def:b1:nernst:couple}{oxydants} plus doux.
\textbf{7.} \ce{C6H12O + 3H2O -> C6H10O4 + 8H+ + 8e-}.
\textbf{8.} $+$V et $+$I.
\textbf{9.} \ce{2HNO3 + 8H+ + 8e- -> N2O + 5H2O}.
\textbf{10.} En additionnant, huit électrons, \ce{8H+} et trois
molécules d'eau se simplifient~:
\ce{C6H12O + 2HNO3 -> C6H10O4 + N2O + 2H2O}.
\textbf{11.} 8.
\textbf{12.} $10^6/146.0 = \qty{6.85e3}{mol}$ d'acide~; $2 \times 6.85 \times
10^3 \times 63.0 = \qty{8.6e5}{g}$, soit \qty{0.86}{t} d'acide nitrique.
\textbf{13.} \ce{H2O2 + 2H+ + 2e- -> 2H2O}.
\textbf{14.} \ce{C6H12O + 4H2O2 -> C6H10O4 + 5H2O}.
\textbf{15.} $4 \times 6.85 \times 10^3 \times 34.0 = \qty{9.3e5}{g}$,
soit \qty{0.93}{t}.
\textbf{16.} L'eau~: aucun oxyde d'azote.
\textbf{17.} Une grande constante d'équilibre ne dit rien de la
vitesse~; \ce{H2O2} réagit lentement sans \omterm{def:b1:elementary-steps:catalyst}{catalyseur}, comme le montre sa
propre décomposition.
\textbf{18.} Il réduit la cyclohexanone en cyclohexanol~:
\ce{4C6H10O + BH4- -> B(OC6H11)4-}, puis
\ce{B(OC6H11)4- + 4H2O -> 4C6H11OH + B(OH)4-}.
\textbf{19.} \qty{6.85e3}{mol}.
\textbf{20.} \qty{6.85e3}{mol} de \ce{N2O} (une mole par mole d'acide).
\textbf{21.} $6.85 \times 10^3/0.95 \times 100.0 = \qty{7.2e5}{g}$, soit \qty{0.72}{t}.
\textbf{22.} $6.85 \times 10^3 \times 44.0 = \qty{3.0e5}{g}$~: \textbf{\qty{0.30}{t}
de \ce{N2O} par tonne d'acide adipique}.
\end{solution}
